\scanpage{52}
\mspage{8}
% archive sheet number: 51
On a ainsi la notion de \uline{graphe pondéré} \deleted{\unsure{rigidifié}} \inserted{\uline{à repères}}
\margin{$\unsure{r|\Gamma}$}
\diagram{
\begin{tikzcd}[ampersand replacement=\&]
\vec R_{\ell}^{+}\ \deleted{\unsure{D_{\infty}}}\subset
\underset{\sigma_0,\sigma_1}{\vec R}\arrow[r] \&
\vec A\arrow[loop right,"\sigma"]\arrow[d,swap,"o"]\&\\
\& S\arrow[r,"\gamma"]\&\mathbb N
\end{tikzcd}
}
Groupoïde (\uline{Gr.pondrep})

foncteur
\diagram{
\begin{tikzcd}[ampersand replacement=\&]
S_{*}\unsure{\mathscr T|\mathcal T}D_{\mathrm{top}}^{\pm}\arrow[r]\&(\mathrm{Gr\ pond\ rep})
\end{tikzcd}
}
qui sera importante sans doute dans l'explicitation de l'opération de ``recollement''. Il est clair en tous cas que toute \unsure{description-traduction} algébrico-combinatoire de $S\mathscr T D_{\mathrm{top}}^{\pm}$ devra permettre, pour chaque objet des \deleted{la catégorie}\inserted{groupoïdes-}paradigmes, de reconstituer (à isom. unique près) le \uline{graphe pondéré à repères} correspondant.

--- plus précisément, il faut reconstituer le foncteur canonique
\diagram{
\begin{tikzcd}[ampersand replacement=\&]
S_{*}\mathscr T D_{\mathrm{top}}^{\pm}\arrow[r]\&(\mathrm{Gr.pondrep.})
\end{tikzcd}
}
(et non seulement le foncteur à valeurs dans (Grbipond))

\uline{Opérations essentielles}

A) Opérations ``d'oubli'' ou de ``\uline{gommage}''.

\textcircled{1} \uline{Gommage de repères} À un objet $(X,\Sigma,R)$, muni en plus d'une partie $R'\subset R$, il associe $(X,\Sigma,R\setminus R')$ (gommage de $R'$).
